\section{At least three vertices with unprotected far targets}
\label{sec:three-unprotected-vertices}

The following argument uses protection rather than the numerical
residual of vertices outside the two candidate sources. In particular,
it is stronger than excluding a two-vertex positive residual support.
All far sets, convenience conditions, and protected pairs refer to
the original oriented graph $D$.

\begin{theorem}\label{thm:three-unprotected-vertices}
Let $D$ be a finite nonempty strongly connected all-deficient oriented graph minimal
under deletion of arbitrary nonempty arc sets. Then at least three
vertices have an unprotected far target:
\[
 \bigl|\{x:t(x)>0\}\bigr|\ge3.
\]
\end{theorem}
\begin{proof}
A directional failure witness of a bad missing pair always has an
unprotected far target: a witness in the opposite direction would,
under protection, create a forbidden two-step path. Thus, if at most
one vertex had an unprotected far target, every missing pair would
be good. Ghazal's good-pair theorem would give a Seymour vertex
\cite[Theorem~2]{Ghazal2012}.

Suppose therefore that the only vertices with $t(x)>0$ are $u,v$.
There is a bad pair, again by the good-pair theorem, and all its
directional failure sources lie in $\{u,v\}$.
Lemma~\ref{lem:two-source-bipartite-bad} gives the missing pair $uv$,
goodness of every missing pair incident to either vertex, and the
nonempty sets
\[
 \mathcal A=N_D^+(u)\cap U_D(v),\qquad
 \mathcal B=N_D^+(v)\cap U_D(u).
\]
The bad pairs are exactly $\mathcal A\times\mathcal B$.
Orient all of them as $\mathcal A\to\mathcal B$, so their only
directional failure source is $u$.

We will canonically orient all good pairs using a reference linear
extension in which every original missing partner of $v$ precedes $v$.
Such an extension exists. Missing partners are incomparable with $v$
in the protected partial order: a protected chain would force an arc.
Adding all precedence requirements $r\prec v$ for
$r\in\mathcal M(v)$ cannot create a cycle, because such a cycle would
require an original protected chain from $v$ to one of these partners.
Let $T$ be the resulting tournament.
By Lemma~\ref{lem:uniform-zero-protected-completion}, every vertex
outside $\{u,v\}$ remains deficient in this one tournament. It is
enough to prove that $v$ also remains deficient.

\par\medskip\noindent\textbf{The empty far-in set.}\quad 
If $P(v)=\varnothing$, every other vertex reaches $v$ within two
steps. Thus $r\to v$ is convenient for every missing partner $r$.
The chosen reference extension makes every added incident arc point
into $v$, whether the reverse direction is also convenient or not.
No first neighbour is added at $v$. Its original first neighbours
cannot introduce a new far second head: good directions are convenient,
and the selected bad directions fail only at $u$.
Hence $g_T(v)=g_D(v)>0$.

Assume henceforth that $P=P(v)\ne\varnothing$ and put
\[
 \begin{gathered}
 Q=N_D^{--}(v),\quad C=\partial P,\quad X=\partial^2P,\quad H=Q\setminus C,\\
 B=H\cap\mathcal M(v),\quad k=|B|,\quad s=|H\cap N_D^+(v)|.
 \end{gathered}
\]
Delete the $s$ outgoing arcs from $v$ into $H$, obtaining
$\widehat D$; when $s=0$ leave $D$ unchanged. Define
\[
 E=(\mathcal M(v)\setminus B)\,\dot\cup\,X,
 \qquad L=E\setminus N^{++}_{\widehat D}(v),
 \qquad p=|P|.
\]
Lemma~\ref{lem:partner-head-budget} gives
\begin{equation}\label{eq:two-support-lambda}
 |L|\le\lambda:=
 \begin{cases}
  \eta(v)-2k,&s=0,\\
  \eta(v)-g_D(v)-2k,&s>0.
 \end{cases}
\end{equation}
It also gives convenience of $r\to v$ for every $r\in B$.

\par\medskip\noindent\textbf{Which missing partners become first neighbours.}\quad 
Write $K=N_T^+(v)\setminus N_D^+(v)$ and $h=|K|$.
Our reference order selects $r\to v$ whenever that direction is
convenient. Consequently every $r\in K$ has an original predecessor
$a\in P(v)$ with $a\to r$: otherwise $r\to v$ would be convenient.
Every exceptional partner in $B$ is directed into $v$, so
$K\subseteq\mathcal M(v)\setminus B\subseteq E$.

For a missing partner $r\notin B\cup\{u\}$, one has $t(r)=0$.
It cannot belong to $P(v)$, because its far target $v$ would then be
protected, forcing the original arc $v\to r$ despite the missing pair.
Thus $r\in C$, and some $a\in P$ points to it. The direction
$r\to v$ is not convenient, so the good pair $rv$ is canonically
oriented as $v\to r$. The only nonexceptional partner which might
instead be directed into $v$ is $u$. Therefore
\begin{equation}\label{eq:two-support-first-count}
 h\ge\mu(v)-k-1.
\end{equation}

\par\medskip\noindent\textbf{A shared envelope for all added first intermediates.}\quad 
We claim
\begin{equation}\label{eq:two-support-shared-envelope}
 N_T^{++}(v)\subseteq
 \bigl(N_D^{++}(v)\cup L\bigr)\setminus K,
 \qquad K\subseteq N_D^{++}(v)\cup L.
\end{equation}
Original first intermediates do not introduce far heads, because the
chosen bad directions fail only at $u$ and good directions are
convenient. Consider an added first intermediate $r\in K$ and an
original far target $y\in U_D(v)\setminus L$.
The original partner-budget lemma already rules out $r\to y$ in $D$.

If $ry$ was missing, it is good: $v$ originally points to neither $r$
nor $y$, so it cannot witness a failure of either direction; the sole
remaining possible source $u$ cannot supply both opposite witnesses.
Choose $a\in P$ with $a\to r$, as above.
Every path of length at most two from $P$ stays inside $P\cup C\cup X$.
Moreover every vertex of this union which is originally far from $v$
lies in $L$. Indeed, such a vertex is either a missing partner outside
$B$ or a member of $X$, hence is in $E$; an originally far target cannot
become reachable after deleting arcs and so is not in
$N^{++}_{\widehat D}(v)$.
Thus $y\notin P\cup C\cup X$ and $y\in U_D(a)$.
The direction $r\to y$ is not convenient, so the good pair $ry$ is
oriented uniquely as $y\to r$. This excludes a new path through $r$
to $y$ and proves the first inclusion.

For the second inclusion, $K\subseteq E$. If a member of $K$ belongs
to $N^{++}_{\widehat D}(v)$, its two-step path already existed in $D$;
it was not an original first neighbour, since it was missing from $v$.
It therefore belongs to $N_D^{++}(v)$. Otherwise it belongs to $L$.
Notice that $L$ may also contain original second neighbours lost in
$\widehat D$, and may contain deleted original first neighbours.
No disjointness between $L$ and $N_D^{++}(v)$ is assumed.

Removing the whole set $K$ from the union in
\eqref{eq:two-support-shared-envelope} gives
\begin{equation}\label{eq:two-support-gap}
 d_T^{++}(v)\le d_D^{++}(v)+\lambda-h,
 \qquad g_T(v)\ge g_D(v)+2h-\lambda.
\end{equation}
If $h=0$, its first neighbourhood is entirely original and has no
chosen-direction failures, so $g_T(v)=g_D(v)>0$ independently of this
upper estimate. We can thus assume $h>0$.

\par\medskip\noindent\textbf{The nonempty deleted bundle.}\quad 
If $s>0$, substitute \eqref{eq:two-support-first-count} and
\eqref{eq:two-support-lambda} into \eqref{eq:two-support-gap}:
\[
 g_T(v)\ge2g_D(v)+2\mu(v)-2-\eta(v)
       =p+g_D(v)-1\ge1.
\]
The equality uses the original identity
$p=2\mu(v)+g_D(v)-1-\eta(v)$.
Only the original graph's arc-set minimality is used to obtain the
nonempty-bundle bound in \eqref{eq:two-support-lambda}.

\par\medskip\noindent\textbf{No deleted outgoing arc and the singleton exception.}\quad 
If $s=0$, the same substitution gives $g_T(v)\ge p-1$.
This is positive when $p\ge2$.
For $p=1$ it is positive unless
$h=\mu(v)-k-1$: any larger integer $h$ improves the estimate by two.
In this last case $u\notin B$ and the canonical direction on $uv$
is $u\to v$.
Write $P=\{a\}$.
If $a\ne u$, then $u$ is a nonexceptional missing partner outside
$P$, hence belongs to $C$. Thus $a\to u$, while $v\in U_D(a)$.
This makes $u\to v$ nonconvenient, contradicting its canonical choice.
If $a=u$, then $C=N_D^+(u)$. Since $s=0$, all of $H=B$ consists of
missing partners of $v$, as does the singleton $P=\{u\}$.
The partition into $P,C,H$ therefore gives
$N_D^+(v)\subseteq C=N_D^+(u)$.
But then
$\mathcal B=N_D^+(v)\cap U_D(u)=\varnothing$,
contrary to the nonempty bad-pair bipartition.
This excludes the only zero-gap estimate case.

\par\medskip\noindent\textbf{Two Seymour vertices in one fixed tournament.}\quad 
We have proved $g_T(v)>0$ for the selected completion.
Every other vertex except possibly $u$ has $t(z)=0$ and remains
deficient by Lemma~\ref{lem:uniform-zero-protected-completion}.
Thus $T$ has at most one Seymour vertex.
The original all-deficient graph has positive out-degree at every
vertex, and adding arcs cannot create a sink. Havet--Thomasse's
two-Seymour-vertex theorem for tournaments with no sink
\cite[Theorem~2]{HavetThomasse2000} is a contradiction.
The comparison is in this one tournament, not between completions
reoriented separately for different candidate feed vertices.
\end{proof}

\begin{corollary}\label{cor:three-positive-residual-support}
Every finite nonempty strongly connected arc-minimal all-deficient oriented graph has
at least three vertices with positive residual:
\[
 \bigl|\{x:\eta(x)>0\}\bigr|\ge3.
\]
In particular, a two-vertex positive residual support is impossible,
without any upper bound on either residual.
\end{corollary}
\begin{proof}
Zero residual makes every far target protected, so $\eta(x)=0$
implies $t(x)=0$. The preceding theorem therefore applies.
\end{proof}

\begin{remark}
The theorem bounds the number of vertices with an unprotected far
target, not the number which actually occur as directional failure
sources of bad missing pairs. A third vertex with $t>0$ need not be
a bad-pair witness; no claim that there must be three actual bad
failure sources is made here.
\end{remark}
